% Uniform vertical gap between major blocks (Education + Experience outer lists, Education subsections)
\def\ResumeSectionBlockSep{10pt}
% Left margin of outer experience/education itemize (body text starts this far in; \section titles do not).
\def\ResumeOuterListLeftMargin{0.15in}
% topsep=0: first job uses the same skip as between jobs (see \vspace after \section{Experience} in experience.tex).
\newcommand{\resumeSubHeadingListStart}{%
  \begin{itemize}[leftmargin=\ResumeOuterListLeftMargin, label={}, itemsep=\ResumeSectionBlockSep, topsep=0pt, parsep=0pt]%
}
\newcommand{\resumeSubHeadingListEnd}{\end{itemize}}

\newcommand{\resumeEduOuterListStart}{%
  \begin{itemize}[leftmargin=\ResumeOuterListLeftMargin, label={}, itemsep=0pt, topsep=0pt, parsep=0pt]%
}
\newcommand{\resumeEduOuterListEnd}{\end{itemize}}

\newcommand{\resumeSubheading}[4]{%
  \item%
  \begin{tabular*}{\textwidth}[t]{l@{\extracolsep{\fill}}r}%
    \textbf{#1} & #2 \\%
    \textit{\small #3} & \textit{\small #4} \\%
  \end{tabular*}\vspace{-5pt}%
}

% Shared bullet list layout (same marker + indent for all resume bullets)
\newcommand{\ResumeBulletLabel}{\textbullet}

% Experience header: #1 role, #2 company (main line); #3 dates | #4 location (subtitle; place in italics).
% Top skip: \ResumeBlockHeadingTopSkip in cv.tex.
\newcommand{\resumeEntryHeading}[4]{%
  \ifhmode\par\fi\addvspace{\ResumeBlockHeadingTopSkip}%
  \item[]%
  {\large\bfseries\rmfamily #1}\ {\normalsize\rmfamily #2}\\*[-0.5ex]%
  {\normalsize\rmfamily #3~\textbar~\textit{#4}}%
}
\let\resumeExperienceHeading\resumeEntryHeading

% Education header: degree + school (single line; no location).
% Top skip: \ResumeBlockHeadingTopSkip in cv.tex.
\newcommand{\resumeEducationHeading}[2]{%
  \ifhmode\par\fi\addvspace{\ResumeBlockHeadingTopSkip}%
  \item[]%
  {\large\bfseries\rmfamily #1}\ {\normalsize\rmfamily #2}%
}

% Degree + institution under the name block (heading.tex): same typography as \resumeEducationHeading but not a list item.
\newcommand{\resumeHeadingDegree}[2]{%
  \par\addvspace{\ResumeHeadingDegreeTopSkip}%
  \noindent{\large\bfseries\rmfamily #1}\ {\normalsize\rmfamily #2}%
}

% In-education subsection: same format as degree/school lines — #1 \large\bfseries, #2 \normalsize on one line (omit #2 if empty).
\newcommand{\resumeEduSubsection}[2]{%
  \par
  \noindent{\large\bfseries\rmfamily #1}%
  \if\relax\detokenize{#2}\relax
  \else
    \ {\normalsize\rmfamily #2}%
  \fi
  \par
}
% Kept as a no-op for older templates; spacing between education blocks is controlled by \resumeEduOuterListStart.
\newcommand{\resumeEduSubsectionBlockGlue}{}

\newcommand{\resumeItemListStart}{%
  \begin{itemize}[leftmargin=0.26in, label=\ResumeBulletLabel, labelsep=0.42em, itemsep=2pt, parsep=0pt, topsep=4pt, partopsep=0pt, after=\vspace{-6pt}]%
}
\newcommand{\resumeItemListEnd}{\end{itemize}}

% Full-width note inside \resumeItemListStart: left flush with \section title; width \textwidth so the right edge matches job headings.
\newcommand{\resumeItemListNoteNoIndent}[1]{%
  \item[]%
  \leavevmode
  \hspace*{\dimexpr-\leftmargin-\labelwidth-\labelsep-\ResumeOuterListLeftMargin\relax}%
  \begin{minipage}[t]{\textwidth}%
  #1%
  \end{minipage}%
}

% Education bullet lists: no topsep (subsection already adds \ResumeSectionTitleSkip below title); no negative after (keeps subsection spacing even).
\newcommand{\resumeEduItemListStart}{%
  \begin{itemize}[leftmargin=0.26in, label=\ResumeBulletLabel, labelsep=0.42em, itemsep=2pt, parsep=0pt, topsep=0pt, partopsep=0pt]%
}

\newcommand{\resumeItem}[1]{\item\small{#1}}
